\chapter{Creating, Displaying, Accessing, and Modifying Data Frames}
\label{chap:data-frames}

\chapterguide
  {You should already understand vectors, lists, indexing, and basic control flow.}
  {create data frames directly and from matrices/lists; display values and structure; inspect dimensions and summary output; select columns and rows; add columns and rows; and remove rows or columns.}

\section{Creating a data frame}

Lab 6a starts with a small employee table. Each column is created from a vector of equal length:

\begin{lstlisting}[style=dsr]
emp.data <- data.frame(
  emp_id = c(1:5),
  emp_name = c("Rick", "Dan", "Michelle", "Ryan", "Gary"),
  salary = c(623.3, 515.2, 611.0, 729.0, 843.25)
)
\end{lstlisting}

A data frame therefore gives the course a tabular structure in which named columns can hold different kinds of data.

\section{Creating a data frame from another structure}

\subsection{From a matrix}

\begin{lstlisting}[style=dsr]
my_matrix <- matrix(c(1, 2, 3, 4, 5, 6), nrow = 2)
print(my_matrix)

df_from_matrix <- data.frame(my_matrix)
print(df_from_matrix)

names(df_from_matrix) <- c("col_1", "col_2", "col_3")
print(df_from_matrix)
\end{lstlisting}

The exercise converts an existing matrix to a data frame and then replaces the automatically generated column names.

\subsection{From a list of vectors}

\begin{lstlisting}[style=dsr]
my_list <- list(
  rating = 1:4,
  animal = c("koala", "hedgehog", "sloth", "panda"),
  country = c("Australia", "Italy", "Peru", "China"),
  avg_sleep_hours = c(21, 18, 17, 10)
)

print(my_list)
super_sleepers <- data.frame(my_list)
print(super_sleepers)
\end{lstlisting}

This connects Lab 3's list concept with a rectangular data structure.

\section{Displaying and inspecting a data frame}

The handout first prints the employee table:

\begin{lstlisting}[style=dsr]
print(emp.data)
\end{lstlisting}

It also points to RStudio's table-view icon and asks you to identify the Console command produced by that interface. The learning objective is to distinguish textual Console output from a table viewer while recognizing that both display the same data-frame object.

\subsection{Structure}

\begin{lstlisting}[style=dsr]
str(emp.data)
\end{lstlisting}

The structure output is used to inspect columns and their stored types.

\subsection{Summary and dimensions}

The supplied handout presents these commands in adjacent subsections:

\begin{lstlisting}[style=dsr]
print(summary(emp.data))

print(dim(emp.data))
print(ncol(emp.data))
print(nrow(emp.data))
\end{lstlisting}

\texttt{summary()} produces per-column summaries, while \texttt{dim()}, \texttt{ncol()}, and \texttt{nrow()} report the table shape.

\begin{pitfall}
The Lab 6a subsection labels appear swapped: the heading ``Display the number of rows and columns'' is followed by \texttt{summary()}, while ``Display the summary'' is followed by \texttt{dim()}, \texttt{ncol()}, and \texttt{nrow()}. These notes organize the functions by what their code actually requests, while preserving this mismatch as a source note.
\end{pitfall}

\section{Accessing columns and rows}

The salary column is extracted in three forms:

\begin{lstlisting}[style=dsr]
print(emp.data$salary)
print(emp.data[["salary"]])
print(emp.data[3])
\end{lstlisting}

The first two forms directly access the named column contents. The third selects the third column using data-frame indexing.

The first two rows with all columns are selected using:

\begin{lstlisting}[style=dsr]
result <- emp.data[1:2, ]
print(result)
\end{lstlisting}

The same row/column coordinate idea seen in matrices is reused, but the result remains tabular.

\section{Adding columns}

Lab 6a adds a department column by assigning a new named component:

\begin{lstlisting}[style=dsr]
emp.data$dept <- c("IT", "Operations", "IT", "HR", "Finance")
print(emp.data)
\end{lstlisting}

It also demonstrates \texttt{cbind()} to attach a start-date column:

\begin{lstlisting}[style=dsr]
emp.newdata <- cbind(
  emp.data,
  start_date = as.Date(c(
    "2012-01-01",
    "2013-09-23",
    "2014-11-15",
    "2014-05-11",
    "2015-03-27"
  ))
)
\end{lstlisting}

This example is the handout's explicit invitation to store multiple data types in one table, including dates.

\section{Adding rows}

The supplied pattern creates a second data frame with the same structure and then combines it with \texttt{rbind()}:

\begin{lstlisting}[style=dsr]
emp.newdata <- data.frame(
  emp_id = c(6:8),
  emp_name = c("Rasmi", "Pranab", "Tusar"),
  salary = c(578.0, 722.5, 632.8),
  dept = c("IT", "Operations", "Finance")
)

emp.finaldata <- rbind(emp.data, emp.newdata)
print(emp.finaldata)
\end{lstlisting}

The handout explicitly notes that rows being appended should have the same structure as the existing data frame.

\section{Removing rows and columns}

Negative indexing is reused to exclude a row and a column:

\begin{lstlisting}[style=dsr]
# Remove the second row and first column
result <- emp.data[-c(2), -c(1)]
print(result)
\end{lstlisting}

This connects directly to the negative-index vector behavior introduced in Lab 3.

\section{Extra practice table}

Lab 6a ends with a friend table containing \texttt{friend\_id}, \texttt{friend\_name}, and \texttt{location}, then asks students to extend the structure with additional rows, columns, and data types.

\begin{lstlisting}[style=dsr]
friends <- data.frame(
  friend_id = 1:5,
  friend_name = c("Sachin", "Sourav", "Dravid", "Sehwag", "Dhoni"),
  location = c("Kolkata", "Delhi", "Bangalore", "Hyderabad", "Chennai")
)
print(friends)
\end{lstlisting}

\begin{quickcheck}
\begin{enumerate}
  \item What does \texttt{str(emp.data)} help you inspect?
  \item Which functions in the handout report the number of rows and columns?
  \item Give two supplied ways to access the salary column by name.
  \item Which function binds new columns and which function binds new rows?
  \item How is negative indexing reused to remove data-frame rows or columns?
\end{enumerate}
\end{quickcheck}

\begin{chapterrecap}
\begin{itemize}
  \item Data frames can be created directly, from matrices, or from lists of equal-length vectors.
  \item Use \texttt{str()} for structure, \texttt{summary()} for summaries, and \texttt{dim()}/\texttt{nrow()}/\texttt{ncol()} for shape.
  \item Columns can be selected by \texttt{\$}, double brackets, or index position.
  \item Add columns with named assignment or \texttt{cbind()}, and add compatible rows with \texttt{rbind()}.
  \item Negative row/column indices remove selected portions in the handout example.
\end{itemize}
\end{chapterrecap}
